\documentclass[11pt]{article}
\usepackage[margin=1in]{geometry}
\usepackage[T1]{fontenc}
\usepackage{amsmath,amssymb,amsthm,mathtools}
\usepackage{microtype}
\usepackage{xcolor}
\usepackage[colorlinks=true,linkcolor=blue!55!black,
  citecolor=blue!55!black,urlcolor=blue!55!black]{hyperref}
\usepackage[nameinlink,noabbrev]{cleveref}
\usepackage{enumitem}
\usepackage{aliascnt}

\newtheorem{theorem}{Theorem}[section]
\newaliascnt{lemma}{theorem}
\newtheorem{lemma}[lemma]{Lemma}
\aliascntresetthe{lemma}
\newaliascnt{proposition}{theorem}
\newtheorem{proposition}[proposition]{Proposition}
\aliascntresetthe{proposition}
\newaliascnt{corollary}{theorem}
\newtheorem{corollary}[corollary]{Corollary}
\aliascntresetthe{corollary}
\newaliascnt{definition}{theorem}
\newtheorem{definition}[definition]{Definition}
\aliascntresetthe{definition}
\theoremstyle{remark}
\newaliascnt{remark}{theorem}
\newtheorem{remark}[remark]{Remark}
\aliascntresetthe{remark}

\newcommand{\LCS}{\operatorname{LCS}}
\newcommand{\rev}{\operatorname{rev}}
\newcommand{\ones}{\#_1}
\newcommand{\zeros}{\#_0}
\newcommand{\Trim}{\operatorname{Trim}}
\newcommand{\polylog}{\operatorname{polylog}}
\newcommand{\TV}{\operatorname{TV}}
\newcommand{\eps}{\varepsilon}
\newcommand{\ind}{\mathbf 1}
\newcommand{\BestMatch}{\operatorname{BestMatch}}
\newcommand{\Match}{\operatorname{Match}}
\newcommand{\Greedy}{\operatorname{Greedy}}

\title{A near-linear-time $(2-c)$-approximation for binary LCS\\
via string regularity}
\author{}
\date{July 13, 2026}

\begin{document}
\maketitle

\begin{abstract}
We give absolute constants $\delta,c>0$ and a deterministic algorithm
which, on arbitrary binary strings $s,t$ of total length
$n=|s|+|t|$, runs in $n\polylog n$ time and outputs an explicit common
subsequence of length at least
$(1/2+\delta)\LCS(s,t)$.  Equivalently, the algorithm is a
$(2-c)$-approximation, with no assumption that the input lengths or symbol
counts agree.  He and Li previously showed that, for every fixed
$\varepsilon>0$, some gain $\delta_\varepsilon>0$ over $1/2$ is attainable
in deterministic $n^{1+\varepsilon}$ time for the same regime.  Thus the
present result achieves a single absolute gain, independent of the exponent
slack $\varepsilon$, in $n\polylog n$ time while retaining explicit-subsequence
output.  Combined with the reduction of Akmal and
Vassilevska Williams,
the result also gives, for every fixed $q\ge2$, a deterministic near-linear
$(1/q+\delta_q)$-approximation with an explicit witness for arbitrary-length
$q$-ary LCS.

The technical core is a promise algorithm for equal-length strings having
exactly half zeros and half ones: it always finds $N/2$ symbols and beats
$N/2$ by a constant fraction when the two strings have a near-perfect common
subsequence.  Its structural input is the string-regularity decomposition of
Guruswami--He--Li; physical-position grids, robust flag-transfer lemmas, and
an exact diagonal-total-variation dynamic program turn the decomposition into
an explicit near-linear-time matching.  A trimming lemma first extends this
core to nearly balanced pairs.  We then give a robust version of the
Rubinstein--Song equal-length case analysis, replacing its edit-distance call
by the trimming lemma, and finally apply the He--Li reduction from equal to
arbitrary input lengths.
\end{abstract}

\begin{quote}
\small\noindent\textbf{Note.} The balanced proof underlying this paper was
fully AI generated by GPT 5.6 Sol, prompted by Yang Liu.  The reduction to
arbitrary inputs and the present revisions were generated and audited by
GPT-5 in the Codex harness, also prompted by Yang Liu.
\end{quote}

\section{Introduction}

Throughout, strings are binary.  For a string $x$, let $\ones(x)$ and
$\zeros(x)$ denote its numbers of ones and zeros.  A common subsequence is
always represented by a noncrossing matching, so every lower bound constructed
below is algorithmically certifiable.

For strings $x$ and $y$, insertion--deletion edit distance and LCS are linked
by the exact identity
\[
 \operatorname{ED}_{\rm ins/del}(x,y)
 =|x|+|y|-2\LCS(x,y).
\]
The standard dynamic program computes either quantity in
$O(|x||y|)$ time.  For strings of lengths $p\ge q$ over a fixed alphabet, the
Four Russians method of Masek and Paterson computes an optimal alignment in
$O(p\max\{1,q/\log p\})$ time on their logarithmic-cost RAM
model~\cite{MasekPaterson}.  Counting its $O(\log n)$-bit boundary vectors as
unit-cost words gives the familiar $O(n^2/\log^2 n)$ equal-length operation
count.  For general sequences, later tabulation refinements compute the LCS
length of strings of lengths $m$ and $n$ in
$O(mn\log\log n/\log^2 n)$ time on a word RAM~\cite{Grabowski}.  These
improvements do not change the quadratic exponent.  Under the Strong
Exponential Time Hypothesis, LCS on constant-size alphabets has no
$O(n^{2-\varepsilon})$ algorithm~\cite{AbboudBackursWilliams}, and this
remains true already for binary strings~\cite{BringmannKunnemann}.

\paragraph{Approximation landscape.}
We use the maximization convention that an $\alpha$-approximation outputs a
common subsequence of length at least $\alpha\LCS(x,y)$; the conventional
approximation factor is $1/\alpha$.  On an alphabet of size $q$, the longest
unary common subsequence gives a $1/q$-approximation in linear time.  For two
length-$n$ strings over an unrestricted alphabet, Hajiaghayi, Seddighin,
Seddighin, and Sun gave a randomized linear-time
$O(n^{0.497956})$-factor approximation in expectation~\cite{HSSS}.
Bringmann, Cohen-Addad, and Das improved this to a
$\widetilde O(n^{0.4})$ factor in linear time with high
probability~\cite{BringmannCohenAddadDas}, and Nosatzki obtained a randomized
$O(n)$-time $n^{o(1)}$-factor approximation with high
probability~\cite{NosatzkiLCS}.  On the deterministic side, Boneh, Golan, and
Kraus gave the first near-linear-time sublinear-factor guarantee, an
$O(n^{3/4}\log n)$ factor, again for two length-$n$
sequences~\cite{BonehGolanKraus}.  These large-alphabet results address a
different regime from the constant gain over $1/2$ sought here for binary
strings.  Hajiaghayi et al. explicitly return a common subsequence;
Bringmann et al. state their main guarantee for the value and note that a
witness can be recovered; Nosatzki's result is a value estimate.

Approximate edit distance developed along a more favorable trajectory.  The
work cited in this paragraph permits substitutions; its metric differs from
insertion--deletion distance by at most a factor of two.
Randomized algorithms of Andoni and Onak achieved a
$2^{\widetilde O(\sqrt{\log n})}$ approximation in $n^{1+o(1)}$
time~\cite{AndoniOnak}, and Andoni, Krauthgamer, and Onak obtained a
$(\log n)^{O(1/\varepsilon)}$ approximation in $n^{1+\varepsilon}$
time~\cite{AndoniKrauthgamerOnak}.  Chakraborty, Das, Goldenberg, Kouck\'y,
and Saks gave the first randomized constant-factor approximation in
$\widetilde O(n^{12/7})$ time~\cite{CDGKS}; Andoni and Nosatzki then achieved
a randomized constant-factor approximation in $n^{1+\varepsilon}$ time for
every fixed $\varepsilon>0$~\cite{AndoniNosatzki}.

Rubinstein and Song showed that an edit-distance estimate in
$[\operatorname{ED},c\operatorname{ED}+o(n)]$ obtained in time $T(n)$ yields
a $(1/2+\delta(c))$-approximation for equal-length binary LCS in
$T(n)+O(n)$ time~\cite{RubinsteinSong}.  Their main theorem is phrased as a
value estimate; when the edit-distance routine returns a transformation, the
reduction also returns an explicit common subsequence.  Instantiating the
reduction with near-linear edit-distance routines broke the binary $1/2$
barrier in randomized $n^{1+\varepsilon}$ time for equal-length inputs.  The
equal-length hypothesis is essential to that analysis, in which a
near-perfect LCS corresponds to small edit distance.

For every fixed alphabet size $q$, Akmal and Vassilevska Williams obtained an
equal-length $(1/q+\delta_q)$-approximation.  They also showed that any
arbitrary-length binary improvement transfers to arbitrary-length $q$-ary
LCS~\cite{AkmalVW}.  He and Li then handled unequal-length binary
strings: for every fixed $\varepsilon>0$, they gave a deterministic
$n^{1+\varepsilon}$-time $(1/2+\delta_\varepsilon)$-approximation for some
$\delta_\varepsilon>0$~\cite{HeLi}.  More
quantitatively, their black-box reduction turns an equal-length gain $\delta$
and running time $T(n)$ into an arbitrary-length gain $\delta^A$ and running
time $O((n+T(n))\log^A n)$ for an absolute $A$; it also preserves an explicit
subsequence when the oracle returns one.  This is the result most directly
comparable with the present theorem.

At a different time--accuracy tradeoff, Mao and Rubinstein recently obtained,
with high probability, randomized $(1-\varepsilon)$-approximations for LCS on
strings of total length $n$ in
$n^2/2^{(\log n)^{\Omega(1)}}$ time~\cite{MaoRubinstein}.  They also observe
that analogous deterministic time--accuracy tradeoffs would imply major new
circuit lower bounds.  The present paper instead keeps deterministic
near-linear time and seeks only an absolute improvement over the unary $1/2$
baseline.

\paragraph{Present contribution.}
The result below strengthens this runtime--accuracy tradeoff: it obtains a
single absolute gain over $1/2$ in $n\polylog n$ time for arbitrary-length
binary strings, while preserving explicit-subsequence output.  Its new
ingredient is the following balanced promise theorem.
The structural source is the three-type string-regularity theorem of
Guruswami, He, and Li, originally developed to prove limitations for binary
deletion codes~\cite{GHL}.  Here it is applied to a hidden near-perfect common
subsequence and converted into an algorithm by enumerable physical
occurrences, robust transfer lemmas, and exact stitching.

\begin{theorem}[Balanced promise theorem]\label{thm:balanced}
There are absolute constants $\eta,\delta>0$ and a deterministic
$N\polylog N$-time algorithm with the following property.  If
$s,t\in\{0,1\}^N$ satisfy
\[
 \ones(s)=\zeros(s)=\ones(t)=\zeros(t)=N/2,
\]
then the algorithm outputs a common subsequence of $s,t$ of length at least
$N/2$.  If additionally
\[
 \LCS(s,t)\ge (1-\eta)N,
\]
then the output has length at least $(1/2+\delta)N$.
\end{theorem}

The promise theorem is the main technical result.  The reductions in
\Cref{sec:reductions} remove both restrictions on its input.

\begin{theorem}[Arbitrary binary LCS]\label{thm:main}
There are absolute constants $\delta_{\rm all},c>0$ and a deterministic
algorithm which, given arbitrary strings $s,t\in\{0,1\}^*$ and
$n=|s|+|t|$, runs in $n\polylog n$ time and outputs a common subsequence
of length at least
\[
  \left(\frac12+\delta_{\rm all}\right)\LCS(s,t).
\]
Consequently its conventional approximation factor is at most $2-c$.
\end{theorem}

\begin{corollary}[Constant alphabets]\label{cor:q-ary}
For every fixed integer $q\ge2$, there is a constant $\delta_q>0$ and a
deterministic algorithm which, on arbitrary $q$-ary strings $x,y$ of total
length $n$, runs in $n\polylog(n+2)$ time and outputs an explicit common
subsequence of length at least $(1/q+\delta_q)\LCS(x,y)$.
\end{corollary}

\begin{proof}
Apply the arbitrary-length alphabet reduction of Akmal and Vassilevska
Williams~\cite{AkmalVW} to \Cref{thm:main}.  For fixed $q$, their construction
makes $O(q^2)$ calls to the binary algorithm and otherwise takes linear time;
it preserves determinism and lifts the returned binary witnesses to a common
subsequence of the original inputs.  Their counting argument gives, for
example, $\delta_q=2\delta_{\rm all}/q$.  The stated running time follows.
\end{proof}

The binary result is constructive throughout: every local certificate is
stored as a noncrossing matching, trimming retains the original positions,
and the two external reductions used in its proof preserve
explicit-subsequence output.

\paragraph{Proof overview.}
The balanced promise theorem applies the three-type decomposition of
Guruswami, He, and Li to a hidden near-perfect common subsequence.  We
enumerate rounded physical occurrences of its regular blocks, transfer their
flags robustly to the observed inputs, and extract local matching gains in
the Imbalanced, Green, and Blue--Yellow cases.  An exact dynamic program
stitches these gains while charging only the total variation of their
matching diagonals.

To pass to general inputs, we first trim a nearly balanced equal-length pair
to two exactly balanced strings.  Deletions in both strings are charged
separately.  Next we revisit all cases in the Rubinstein--Song reduction.
Their approximate-edit-distance routine occurs only on nearly balanced
equal-length pairs, where the trimming algorithm supplies the required
gain; all cross-unbalanced cases use unary or greedy matchings.  The resulting
equal-length black box is then inserted into the He--Li reduction for
possibly unequal lengths.

\paragraph{Organization.}
\Cref{sec:one-rank} records the regularity input.  The hidden common core,
flag transfer, local matchings, and stitching are developed in
\Cref{sec:hidden-core,sec:flag-transfer,sec:local-modules,sec:stitching}.
\Cref{sec:balanced-algorithm,sec:balanced-completeness} prove the balanced
promise theorem.  \Cref{sec:reductions} gives the nearly balanced,
equal-length, and unequal-length reductions and proves \Cref{thm:main}.

\paragraph{Limitations.}
The result is asymptotic.  The fixed regularity parameters force an enormous
polylogarithmic exponent and an enormous absolute threshold below which the
algorithm falls back to exact dynamic programming.  Neither the approximation
gain nor these quantitative constants are intended to be practical, and the
paper does not improve the near-quadratic regime for near-exact
approximation.

The constants are extremely small.  We make no attempt to optimize them.  We
fix the constants used by Guruswami--He--Li,
\begin{equation}\label{eq:constants}
 \eps=10^{-6},\qquad \gamma=10^{-3}\eps^2=10^{-15},
\end{equation}
and set
\begin{equation}\label{eq:rho-kappa}
 \rho=\eps/100.
\end{equation}
After all absolute constants in the proof are fixed, choose
\begin{equation}\label{eq:hierarchy}
 0<\eta\ll \sigma\ll \kappa\ll \eps^{6}.
\end{equation}
For example, $\kappa=\eps^{20}$, $\sigma=\eps^{30}$, and a sufficiently
small constant multiple of $\eps^{50}$ for $\eta$ are more than adequate.
Only the hierarchy in \eqref{eq:hierarchy} is used.

\section{One-rank notation and the imported regularity theorem}
\label{sec:one-rank}

\subsection{One-rank intervals}

Let $x$ be any string containing $M$ ones.  For
$a\in[M]$ and an integer $b>a$, write $x[a,b)$ for the substring beginning
at the $a$th one and ending immediately before the $b$th one, or at the end
of $x$ when $b>M$.  Thus $x[a,a+r)$ contains exactly $r$ ones when
$a+r\le M+1$ and is truncated otherwise.  We write
\[
 z_x(a,r):=\zeros(x[a,a+r)).
\]
For $a\in[M+1]$, define $x[a,a)$ to be the empty string.
This is the half-open version of the notation in Guruswami--He--Li.

For a positive integer $r$, the start $a$ is an $r$-flag of rate
$z_x(a,r)/(r-1)$, with the usual convention at $r=1$.  It is
\emph{strong Blue}, \emph{strong Green}, \emph{strong Yellow}, or
\emph{strong Red} according as its rate is respectively
\[
 >\eps^{-1},\qquad >1+2\eps,\qquad >0.9,\qquad \le 0.9.
\]
For $a\in[M]$, let $b_x(a)$ be the largest power of two $r$ for which $a$ is
a strong Blue $r$-flag, subject to $r\le M$, and set $b_x(a)=0$ if none
exists.

If $x$ has $B$ ones and $\Delta\in[-B,B]$, define
\begin{equation}\label{eq:trim-def}
 \Trim_\Delta(x)
 :=x\bigl[\max\{1,1-\Delta\},
          \min\{B+1,B+1-\Delta\}\bigr).
\end{equation}
It contains exactly $B-|\Delta|$ ones: it is a prefix when $\Delta>0$
and a suffix when $\Delta<0$.  A start $j$ is \emph{allowable at scale $r$}
in a $B$-one string if $1\le j\le B-r+1$, equivalently if its full
$r$-window is contained in the string.

When $x$ has leading zeros, they are therefore omitted from every one-rank
interval.  The reversal $\rev(x)$ is the reversal convention of
Guruswami--He--Li: keep the initial one fixed and reverse all subsequent
symbols.  It is defined only for strings beginning with one.  It preserves
LCS between two such strings and reverses one-rank intervals exactly.  We
never apply this operation to an arbitrary full input string.

\subsection{The three types}

Let $x$ contain $B=2^m$ ones.  We use the following definitions.

\begin{enumerate}[label=(\roman*)]
\item $x$ is \emph{Imbalanced} if some one-rank interval $I$ of size
      $|I|>\eps^5B$ satisfies
      \[
       \zeros(x_I)\notin (1\pm\eps)|I|.
      \]
\item $x$ is $r$-\emph{Green} if $r\le\eps^5B$ and it has at least
      $\eps^2B$ strong Green $r$-flags.
\item For an integer $m_0\in[1,m]$, $x$ is
      $m_0$-\emph{Blue--Yellow} if at least
      $(\eps^2-\gamma)B$ indices $a$ satisfy
      \[
       2^{m_0}\le b_x(a)\le\gamma B,
      \]
      and, for every integer $r$ with $2^{m_0}<r\le B$, at most
      $600\eps B$ starts are strong Red $r$-flags.
\end{enumerate}

When several types are available, we assign one with priority
\[
 \text{Imbalanced}>\text{Green}>\text{Blue--Yellow}.
\]

For a $B$-one string $x$, call it $6\gamma$-\emph{flag-balanced} if the
$\ell_1$ distance between the distributions of $b_x(a)$ in its first
and second one-rank halves is at most $6\gamma$.  In particular, if $x$ is
both flag-balanced and $m_0$-Blue--Yellow, then its first half contains at
least
\begin{equation}\label{eq:blue-first-half}
 \frac12(\eps^2-4\gamma)B>0.49\eps^2B
\end{equation}
starts $a$ with $2^{m_0}\le b_x(a)\le\gamma B$.

\subsection{Imported theorem}

We use the following consequence of Lemmas 3.5, 4.4, and 5.6 of
Guruswami--He--Li~\cite{GHL}.

\begin{proposition}[GHL regularity package]\label{prop:GHL}
Let $w$ begin with one and contain $L=2^n$ ones, for sufficiently large $n$.
Put
\[
 m_0^{\rm big}=\left\lfloor n-200\gamma^{-3}\log_2 n\right\rfloor,
 \qquad B_0=2^{m_0^{\rm big}}.
\]
At least one of the following alternatives holds.
\begin{enumerate}[label=(\alph*)]
\item In $w$ or $\rev(w)$, more than $\gamma^2L$ indices satisfy
      $b(a)\ge B_0$.  Then that orientation contains pairwise disjoint
      zero-rich Imbalanced one-rank intervals whose total one-rank size is at
      least $\gamma^2L/2$ and each of whose sizes is at least $B_0$.
\item For every fixed constant $R$, there is, for all sufficiently large
      $n$, an integer
      \[
       m_*\in[n-150\gamma^{-3}\log_2 n,\ n-R]
      \]
      such that, with $B=2^{m_*}$ and $K=L/B$, all but $3\gamma K$ of the
      consecutive $B$-one blocks of $w$ are $6\gamma$-flag-balanced, and
      all but $3\gamma K$ of their reversals are
      $6\gamma$-flag-balanced.  Every block and every reversed block has one
      of the three types above.
\end{enumerate}
Moreover $K\le(\log L)^{A_0}$ for an absolute constant $A_0$ depending only
on $\gamma$.
\end{proposition}

\begin{proof}
If more than $\gamma^2L$ indices have $b(a)\ge B_0$, greedily choose the
leftmost such index not covered by a previously selected interval and select
the interval
$[a,\min\{a+b(a)-1,L\}]$, truncated at the end of the string.  These
intervals are disjoint and cover every exceptional index,
so their total one-rank size is at least $\gamma^2L$.  Each is zero-rich
Imbalanced by GHL Lemma 3.5.  Only the final selected interval can be truncated
to size below $B_0$; it then contains fewer than $B_0$ one-ranks.  For
sufficiently large $n$, $B_0<\gamma^2L/2$.  Discarding this possible final
interval leaves total size at least $\gamma^2L/2$, proving (a).

Otherwise at most $\gamma^2L$ indices satisfy $b(a)\ge B_0$.  Since
$B_0\le 2^{n-200\gamma^{-3}\log_2 n}$, the threshold hypothesis of GHL
Lemma 5.6 holds.  That lemma, with $\beta=\gamma$, excludes at most
$32\gamma^{-3}\log_2 n$ scales for $w$ and the same number for $\rev(w)$
inside an interval of $150\gamma^{-3}\log_2 n$ top scales.  After excluding
the fixed top $R$ scales, at least one common scale remains for sufficiently
large $n$.  Lemma 5.6 gives the asserted $6\gamma$-balance at that scale.
GHL Lemma 4.4 supplies a type for every power-of-two-one block.  Finally,
$n-m_*=O(\log n)$ implies
$K=2^{n-m_*}=(\log L)^{O(1)}$.
\end{proof}

\section{The hidden common core and physical-grid candidates}
\label{sec:hidden-core}

Fix balanced strings $s,t\in\{0,1\}^N$ and a common subsequence matching
$q$ of length $N-D$.  Each input has exactly $D$ unmatched symbols.  Since
$s$ has $N/2$ ones and zeros,
\begin{equation}\label{eq:q-counts}
 \ones(q),\zeros(q)\ge N/2-D.
\end{equation}
Let $L=2^n$ be the largest power of two not exceeding $\ones(q)$, and let
$q'=q[1,L+1)$ (ending with $q$ if $q$ has exactly $L$ ones).
Then
\begin{equation}\label{eq:L-linear}
 L>\frac{\ones(q)}2\ge \frac N4-\frac D2.
\end{equation}

At a candidate block scale $B$, partition $q'$ into consecutive $B$-one
blocks $Q_1,\dots,Q_K$.  Let $S_i$ and $T_i$ be the minimal physical spans
of the images of $Q_i$ in $s$ and $t$.  Put
\[
 e_i^s=|S_i|-|Q_i|,\qquad e_i^t=|T_i|-|Q_i|,
 \qquad e_i=e_i^s+e_i^t.
\]

\begin{lemma}[Defect accounting]\label{lem:defect}
The spans $S_i$ are ordered and pairwise disjoint, as are the $T_i$, and
\[
 \sum_i e_i\le 2D.
\]
\end{lemma}

\begin{proof}
The blocks $Q_i$ are consecutive pieces of one fixed noncrossing matching.
No symbol matched to a different $Q_j$ lies inside the minimal span of
$Q_i$.  Hence every symbol of $S_i$ outside the image of $Q_i$ is unmatched
in the fixed matching, and the same is true in $t$.  Disjointness charges
every unmatched symbol at most once.
\end{proof}

Fix $h=\lceil\sigma B\rceil$ and let the physical grid consist of
$0,h,2h,\dots,N$, together with $N$ if necessary.  Round a physical interval
outward to the nearest grid points.

\begin{lemma}[Physical-grid lemma]\label{lem:grid}
Let $\widehat S_i,\widehat T_i$ be the outward rounded spans.  Then they
contain the two images of $Q_i$ and
\[
 (|\widehat S_i|-|Q_i|)+(|\widehat T_i|-|Q_i|)
 \le e_i+4h.
\]
One of the two index parities contains at least half of any prescribed set of
indices and, on either fixed parity, the rounded spans are pairwise disjoint
in each input, provided $2h<B$.
\end{lemma}

\begin{proof}
Outward rounding adds fewer than $2h$ physical symbols per string.  Between
$S_i$ and $S_{i+2}$ lies the entire span $S_{i+1}$, which contains $B$
matched ones and therefore has physical length at least $B$.  Rounding can
reduce this gap by less than $2h$, so it remains positive.  The same argument
applies in $t$.
\end{proof}

The algorithm enumerates every physical grid interval.  At the successful
regular scale, the number of grid points is
\begin{equation}\label{eq:grid-count}
 G=O(N/h)=O(N/B)=\polylog N,
\end{equation}
so there are only $O(G^2)$ intervals per input and $O(G^4)$ interval pairs.
We shall repeatedly use
\begin{equation}\label{eq:Kh}
 Kh\le \sigma L+K=O(\sigma L).
\end{equation}
The last equality holds once $B\ge 2/\sigma$; the finitely many smaller
instances are included in the exact-algorithm fallback below.

If a physical interval $X$ contains at least $B$ ones, let $F_B(X)$ be its
substring beginning at its first one and ending immediately before its
$(B+1)$st one, or at the end of $X$ if there is no such one.  Thus $F_B(X)$
has exactly $B$ ones.  Define $L_B(X)$ analogously from the last $B$ ones,
and use $\rev(L_B(X))$ for reverse-oriented certificates.

\section{Strong-to-weak flag transfer}
\label{sec:flag-transfer}

Define, for every integer $r\ge1$,
\begin{equation}\label{eq:H}
 H(r)=r+\lfloor\rho(r-1)\rfloor.
\end{equation}
The floor is essential: it gives $H(2)=2$ and, more generally,
\begin{equation}\label{eq:H-bound}
 H(r)\le(1+\rho)r.
\end{equation}

\subsection{A maximal inequality}

\begin{lemma}[One-sided maximal inequality]\label{lem:maximal}
Let $f_1,\dots,f_M\ge0$, with $\sum_jf_j=E$.  For $\lambda>0$, the set of
indices $i$ for which some $r\ge1$ satisfies
\[
 \sum_{j=i}^{i+r-1}f_j>\lambda r
\]
has size at most $3E/\lambda$.
\end{lemma}

\begin{proof}
For every bad $i$, choose an interval $I_i=[i,i+r_i)$ witnessing the
inequality.  Order these intervals by nonincreasing length and greedily retain
an interval when it is disjoint from those already retained.  This gives a
disjoint subfamily $\mathcal J$.  Every discarded interval intersects a
retained interval of at least its length, so its left endpoint lies in the
threefold enlargement of that retained interval.  Therefore the number of bad
integer left endpoints is at most $3\sum_{J\in\mathcal J}|J|$.  Since the
chosen intervals are disjoint,
\[
 \lambda\sum_{J\in\mathcal J}|J|
 <\sum_{J\in\mathcal J}\sum_{j\in J}f_j\le E.
\]
This proves the claim.
\end{proof}

When a core string $Q$ is embedded in a supersequence $X$, put into $f_j$
the number of extra one symbols in the gap following the $j$th core one.
If the $r$-one core interval starting at $i$ contains at most
$\lfloor\rho(r-1)\rfloor$ extra ones, the $H(r)$-one interval of $X$ starting
at the image of $i$ contains the entire core interval.  By
\Cref{lem:maximal}, simultaneously over variable $r\ge2$, all but
$O(E/\rho)$ starts have this property; explicitly, the displayed maximal
inequality gives at most $6E/\rho$ after using $r-1\ge r/2$.

\subsection{Canonical-window transfer}

The next lemma records precisely what survives after taking the first $B$
ones of a rounded physical span.  It does \emph{not} claim that this canonical
window is close in ordinary edit distance to the whole core block.

\begin{lemma}[Canonical flag stability]\label{lem:canonical-stability}
Let $Q$ begin with one and contain exactly $B$ ones.  Let $X$ be a physical
string containing $Q$ as a subsequence and at most $a$ additional ones, and
put $U=F_B(X)$.  Assume $a\le\kappa B$.
\begin{enumerate}[label=(\alph*)]
\item If $Q$ has at least $\eps^2B$ strong Green $r$-flags, where
      $r\le\eps^5B$, then there is a scale $r'$ depending only on $r$ such
      that $U$ has at least $0.8\eps^2B$ weak Green $r'$-flags of rate
      greater than $1+1.5\eps$.  The same $r'$ works for every occurrence
      $X$ satisfying the hypotheses.
\item Suppose $Q$ is $6\gamma$-flag-balanced and $m_0$-Blue--Yellow.  Then
      $U$ has at least $0.45\eps^2B$ starts in its first one-rank half, each
      equipped with a power
      $p\in[2^{m_0},\gamma B]$, such that
      \begin{equation}\label{eq:weak-blue}
       z_U(i,H(p))\ge 0.45p/\eps.
      \end{equation}
      Moreover, for every power $p\in[2^{m_0},\gamma B]$, at most
      $700\eps B$ allowable
      starts $j$ fail
      \begin{equation}\label{eq:weak-yellow}
       z_U(j,Q(p))\ge0.46p/\eps,
       \qquad
       Q(p):=H\!\left(\left\lceil0.55p/\eps\right\rceil\right).
      \end{equation}
\end{enumerate}
There is also a right-to-left version: if $\rev(Q)$ has either property,
then the corresponding conclusion holds, in right-to-left one-ranks, in
$\rev(L_B(X))$.
\end{lemma}

\begin{proof}
Write $c_1,\ldots,c_B$ for the one symbols of the displayed copy of $Q$ in
$X$, and let $\phi(k)$ be the one-rank of $c_k$ in $X$.  Since at most $a$
ones of $X$ are outside this copy,
\begin{equation}\label{eq:phi}
 k\le\phi(k)\le k+a.
\end{equation}
In particular $c_1,\ldots,c_{B-a}$ occur in $U$, and at most $a$ mapped
starts cross any fixed one-rank boundary.

For $1\le j<B$, let $f_j$ count the additional ones of $X$ strictly between
$c_j$ and $c_{j+1}$, and let $f_B$ count the additional ones after $c_B$.
Consider first a start $k\le B-r+1$, so that its full core window is
defined.  If
\begin{equation}\label{eq:local-insertions}
 \sum_{j=k}^{k+r-1}f_j\le\lfloor\rho(r-1)\rfloor,
\end{equation}
then the $H(r)$-one window of $X$ starting at rank $\phi(k)$ contains the
entire $r$-one core window starting at $c_k$, including its terminal
one-gap.  For $r\ge2$, failure of \eqref{eq:local-insertions} implies
$\sum_{j=k}^{k+r-1}f_j>(\rho/2)r$.  Hence
\Cref{lem:maximal} shows that, even when $r$ varies with $k$ among full core
windows, at most
$6a/\rho$ starts fail this containment.  We use the round upper bound
$8a/\rho$ below.

For (a), first discard the at most $r-1$ core starts without a full
$r$-window and assume $r\ge2$.  Outside a further set of at most
$8a/\rho$ starts,
the image $H(r)$-window contains the full strong core window.  It therefore
contains more than $(1+2\eps)(r-1)$ zeros.  Put
\[
 d=H(r)-1
\]
and round $d$ upward to a member $g$ of the integer geometric grid from
\Cref{lem:geometric-green} below.  That grid satisfies
\[
 d\le g\le(1+2\eps/100)d.
\]
Extending the window cannot remove zeros, and
\[
 \frac{(1+2\eps)(r-1)}{g}
 \ge
 \frac{1+2\eps}{(1+\rho)(1+2\eps/100)}
 >1+1.5\eps.
\]
Thus $r'=g+1$ works.  The map $k\mapsto\phi(k)$ is injective.  Including the
initially discarded core-boundary starts, at most $a+r'$ starts are lost at
the terminal boundary of $F_B(X)$.  Since
$r'=O(\eps^5B)+O_\eps(1)$ and $a\le\kappa B\ll\rho\eps^2B$, fewer than
$0.2\eps^2B$ starts are lost.  If $r=1$, a strong Green flag means that the
one-gap contains a zero.  Each inserted one destroys this property for at
most one core one-gap, and the same count follows directly.

For (b), \eqref{eq:blue-first-half} gives more than $0.49\eps^2B$ strong
Blue labeled starts in the core first half.  Apply \Cref{lem:maximal} to the
variable powers $p$.  Outside at most $8a/\rho$ starts, the candidate
$H(p)$-window contains the strong core $p$-window.  Since $p\ge2$,
\[
 \eps^{-1}(p-1)\ge p/(2\eps)>0.45p/\eps.
\]
These core windows are full because their starts are in the first half and
$p\le\gamma B<B/2$.
At most $a$ starts move across the first-half boundary, and the maximum
window length is at most $(1+\rho)\gamma B$.  The hierarchy of constants
therefore leaves at least $0.45\eps^2B$ starts satisfying
\eqref{eq:weak-blue}.

Fix a power $p$ and put
\[
 q_0=\left\lceil0.55p/\eps\right\rceil.
\]
Discard the last $q_0-1$ core starts, which have no full $q_0$-window.  The
Blue--Yellow type says that at most $600\eps B$ of the remaining core starts
are strong
Red $q_0$-flags, because $q_0>2^{m_0}$.  Every other core start has more than
$0.9(q_0-1)$ zeros in its $q_0$-window.  Outside $O(a/\rho)$ insertion-dense
starts, the candidate $H(q_0)=Q(p)$ window contains it.  As $p\ge2$ and
$\eps\le10^{-6}$,
\[
 0.9(q_0-1)
 \ge0.495p/\eps-0.9
 >0.46p/\eps.
\]
For this fixed scale, double counting shows that at most $2a/\rho$ core
starts are insertion-dense.  The candidate starts not arising from retained
core starts, the shifted terminal starts, and the discarded $q_0$-boundary
starts account for at most $2a+Q(p)$ further exceptions (the last group is
covered because $Q(p)\ge q_0$).  Finally
\[
 Q(p)\le0.551p/\eps\le0.000551\eps B.
\]
Thus the total number of failures is at most
\[
 600\eps B+2a/\rho+2a+Q(p)<700\eps B
\]
under the chosen constant hierarchy.

For completeness, the right-to-left assertion is not an appeal to literal
subsequence containment under the modified reversal.  A $B$-one string $y$
has a unique representation
\[
 y=1\,0^{g_1}1\,0^{g_2}\cdots1\,0^{g_B},
 \qquad
 \rev(y)=1\,0^{g_B}1\,0^{g_{B-1}}\cdots1\,0^{g_1}.
\]
Consequently, whenever the displayed indices lie in $[B]$,
\[
 z_{\rev(y)}(i,r)=
 \sum_{j=B-i-r+2}^{B-i+1}g_j.
\]
Put $W=L_B(X)$ and enumerate the ones of $Q$ and $W$ from the right.  If
$\psi(k)$ is the right-rank in $W$ of the $k$th core one from the right,
then $k\le\psi(k)\le k+a$.  For every rank $k\ge2$, the $k$th one of
$\rev(Q)$ corresponds to the $(k-1)$st core one from the right, and its image
in $\rev(W)$ has rank $1+\psi(k-1)\in[k,k+a]$.  The intervening gap sequence
is exactly the original gap sequence read right-to-left, with only the
additional ones of $X$ inserted.  Thus \eqref{eq:local-insertions}, the
maximal inequality, and the boundary counts apply verbatim.  The exceptional
rank $k=1$ is discarded.  This one extra loss is absorbed by the
$a+r'$ and $2a+Q(p)$ bounds for sufficiently large $B$.
\end{proof}

\section{Local matching modules}
\label{sec:local-modules}

\subsection{Imbalanced intervals}

\begin{lemma}[Robust Imbalanced certificate]\label{lem:imbalanced}
Suppose a common core interval has $r$ ones and $z$ zeros, with
$r\ge\eps^5B$, and its two physical occurrence spans acquire at most $k$
additional symbols in total.
\begin{enumerate}[label=(\alph*)]
\item If $z\ge(1+\eps)r$, matching zeros gives, relative to the all-one
      baseline, adjusted reward at least $\eps r-k/2$.
\item If $z\le(1-\eps)r$, matching ones gives, relative to the all-zero
      baseline, adjusted reward at least $\eps r-k/2$.
\end{enumerate}
In particular the applicable reward is at least $\eps r/2$ whenever
$k\le\eps r$.
\end{lemma}

\begin{proof}
In case (a), both occurrences contain all $z$ common zeros, while their total
numbers of ones are at most $r+k_s$ and $r+k_t$, where $k_s+k_t=k$.  Hence
\[
 z-\frac{(r+k_s)+(r+k_t)}2\ge\eps r-k/2.
\]
Case (b) is identical after exchanging the symbols: both occurrences contain
the $r$ core ones, and their total zero counts are at most
$2z+k$.
\end{proof}

The algorithm finds these certificates because it enumerates all physical
grid intervals and stores both unary matchings for every interval pair.
Rounding the two images of a witness outward adds at most $4h$ symbols.

\subsection{Green matching}

\begin{lemma}[Integer geometric scales]\label{lem:geometric-green}
There is a set $\mathcal G_B$ of
$O(\eps^{-1}\log B)$ nonnegative integers such that, for every integer
$d\in[0,B]$, its least upper neighbor $g\in\mathcal G_B$ satisfies
\[
 d\le g\le(1+2\eps/100)d
\]
when $d>0$, while $0\in\mathcal G_B$.
\end{lemma}

\begin{proof}
Include every integer up to $100/\eps$.  Above that threshold, include the
ceilings of successive powers of $1+\eps/100$, stopping only after including
a value at least $B$.  For $d>100/\eps$, the
additive ceiling error is at most $(\eps/100)d$, proving the bound.
\end{proof}

Call a start a weak canonical $r$-Green flag if it has at least one zero when
$r=1$, and if its rate is greater than $1+1.5\eps$ when $r\ge2$.

\begin{lemma}[Weak Green matching]\label{lem:green}
Let $U,V$ each contain exactly $B$ ones.  Suppose that for one common scale
$r$ they each have at least
$\alpha B$ weak canonical $r$-Green starts whose $r$-windows are contained
in the string, where $\alpha=0.8\eps^2$.  For
$\Delta\in[-B,B]$, let $M_r(\Delta)$ be the number of starts that are weak
Green in $U$ and whose $\Delta$-shift is weak Green in $V$, with both full
$r$-windows contained in the corresponding trimmed strings.  Then
\begin{equation}\label{eq:green-pointwise}
 \LCS(\Trim_\Delta U,\Trim_{-\Delta}V)
 \ge B-|\Delta|+\frac\eps2M_r(\Delta),
\end{equation}
and consequently
\begin{equation}\label{eq:green-expectation}
 \mathbb E_\Delta\bigl[
  \LCS(\Trim_\Delta U,\Trim_{-\Delta}V)+|\Delta|
 \bigr]
 \ge B+0.05\eps^5B.
\end{equation}
Here and below the expectation is uniform over the displayed integer range
of $\Delta$.
\end{lemma}

\begin{proof}
Equivalently, $M_r(\Delta)$ counts pairs $(i,j)$ of allowable weak starts
with $j-i=\Delta$.  The two full windows are then in the trims: if
$\Delta\ge0$, $j\ge1+\Delta$ and $i+r-1\le B-\Delta$, and the case
$\Delta<0$ is symmetric.  The correlation identity gives
\[
 \sum_{\Delta=-B}^{B}M_r(\Delta)=|G_U||G_V|,
\]
so $\mathbb E M_r(\Delta)\ge\alpha^2B/3$.  For a fixed $\Delta$, greedily
retain $r$-separated paired starts.  At least $M_r(\Delta)/r$ remain.  If
$r\ge2$, replacing the following $r-1$ matched ones by zeros gains at least
$1.5\eps(r-1)\ge0.75\eps r$ per retained pair.  If $r=1$, the common zero in
the one-gap is inserted without deleting a one match and gives a gain of one.
In either case the gain is at least $(\eps/2)M_r(\Delta)$, proving
\eqref{eq:green-pointwise}.  Averaging and using
$\alpha^2\eps/6>0.1\eps^5$ proves the more conservative
\eqref{eq:green-expectation}.
\end{proof}

All $M_r(\Delta)$ for all canonical scales can be computed deterministically
in $O(\eps^{-1}B\log^2B)$ time by binary cross-correlations.  After a scale
and shift are selected, the greedy matching in the proof is reconstructed in
$O(B)$ time.

\subsection{A slackened Blue--Yellow lemma}

\begin{lemma}[Directional weak Blue--Yellow matching]\label{lem:BY-directional}
Let $U,V$ each have exactly $B$ ones, with $B$ sufficiently large in terms
of $\eps$.  Fix an integer $m_0\in[1,\lfloor\log_2B\rfloor]$, and assume:
\begin{enumerate}[label=(\roman*)]
\item at least $0.45\eps^2B$ starts $i\le B/2$ are each assigned a power
      $p(i)\in[2^{m_0},\gamma B]$ satisfying
      \[
       z_U(i,H(p(i)))\ge0.45p(i)/\eps;
      \]
\item for every power $p\in[2^{m_0},\gamma B]$, at most $700\eps B$
      allowable starts $j$ fail
      \[
       z_V(j,Q(p))\ge0.46p/\eps,
       \qquad Q(p)=H(\lceil0.55p/\eps\rceil).
      \]
\end{enumerate}
Put
\[
 S=\left\lfloor(1/2+0.01\eps)B\right\rfloor,
 \qquad
 T=\left\lfloor(1/2+0.27\eps)B\right\rfloor.
\]
For every integer $\Delta\in[0,\lfloor B/4\rfloor]$, one can explicitly
construct a common subsequence of $U[1,S+1)$ and
$V[1+\Delta,1+\Delta+T)$, of weight $W(\Delta)$, such that
\begin{equation}\label{eq:BY-reward}
 \mathbb E_\Delta\left[W(\Delta)-\frac{S+T}{2}\right]
 \ge0.04\eps B.
\end{equation}
\end{lemma}

\begin{proof}
Greedily choose the leftmost remaining qualifying source start $i_k$, set
$p_k=p(i_k)$, and delete all qualifying starts in
$[i_k,i_k+H(p_k))$.  Continue until
\[
 P:=\sum_kp_k\ge0.44\eps^2B.
\]
This stopping point exists: the disjoint chosen intervals cover all
$0.45\eps^2B$ qualifying starts, while
$H(p)\le(1+\rho)p$.  Since the final power is at most
$\gamma B=0.001\eps^2B$,
\begin{equation}\label{eq:P-range}
 0.44\eps^2B\le P\le0.441\eps^2B.
\end{equation}
Every selected source interval lies before $S$, because its start is at most
$B/2$ and $H(p)\le(1+\rho)\gamma B\ll0.01\eps B$.

Write $h_k=H(p_k)$ and $q_k=Q(p_k)$.  For $p\ge2$,
\begin{equation}\label{eq:Q-bound}
 H(p)\le(1+\rho)p,
 \qquad
 0.55p/\eps\le Q(p)\le0.551p/\eps.
\end{equation}
The upper bound uses $p\ge2$ and $\eps\le10^{-6}$.  Therefore
\begin{equation}\label{eq:Qsum}
 \sum_kq_k\le0.551P/\eps\le0.242991\eps B.
\end{equation}

For a shift $\Delta$, define
\begin{equation}\label{eq:j-k}
 j_k=i_k+\Delta+\sum_{r<k}(q_r-h_r).
\end{equation}
The source intervals are disjoint, and hence
\[
 j_{k+1}-(j_k+q_k)=i_{k+1}-(i_k+h_k)\ge0.
\]
The final target endpoint, including the final one-gap, is at most
\[
 1+\Delta+S+\sum_k(q_k-h_k)
 \le1+\Delta+(1/2+0.252991\eps)B
 <1+\Delta+T.
\]
Thus every target interval is allowable.

Call $k$ good if condition (ii) holds at $j_k$.  As $\Delta$ ranges over at
$\lfloor B/4\rfloor+1\ge B/4$ consecutive integers, $j_k$ does likewise, so
\begin{equation}\label{eq:bad-prob}
 \Pr[k\text{ bad}]\le2801\eps
\end{equation}
for sufficiently large $B$.  Put
$b_k=\lfloor0.44p_k/\eps\rfloor$.  Both source and a good target interval
contain at least $b_k$ zeros.

Partition the source window into the gap before the first selected interval,
the selected intervals, the gaps between them, and the final gap ending at
one-rank $S$.  In the target, use the gap beginning at one-rank $1+\Delta$,
the intervals beginning at the $j_k$, the corresponding intermediate gaps,
and the final gap ending at
$1+\Delta+S+\sum_k(q_k-h_k)$.  Equation \eqref{eq:j-k} makes every pair of
corresponding gaps have equal one counts.  All target pieces lie in the
declared $T$-one window by the preceding endpoint bound.  Match all gap ones;
on a bad flag pair match all $h_k$
source ones to target ones; and on a good pair replace those matches by
$b_k$ zero matches.  This is noncrossing and has weight
\begin{equation}\label{eq:W}
 W(\Delta)=S+\sum_{k\text{ good}}(b_k-h_k).
\end{equation}
Every summand is positive.  If $d$ is the number of selected flags, then
$d\le P/2$, because every $p_k\ge2$.  Consequently
\[
 \sum_kb_k\ge0.44P/\eps-P/2,
 \qquad
 \sum_kh_k\le(1+\rho)P.
\]
Using \eqref{eq:bad-prob} and \eqref{eq:P-range},
\begin{align*}
 \mathbb EW(\Delta)
 &\ge S+(1-2801\eps)
       \left(\frac{0.44}{\eps}-1.5-\rho\right)P\\
 &\ge(1/2+0.19\eps)B
\end{align*}
for sufficiently large $B$.  Since
$(S+T)/2\le(1/2+0.14\eps)B$, the conservative bound
\eqref{eq:BY-reward} follows.
\end{proof}

\begin{lemma}[Directional cancellation]\label{lem:BY-pair}
Let $U_1,U_2,V_1,V_2$ each have $B$ ones.  Suppose
\Cref{lem:BY-directional} applies to $(U_1,V_1)$ and to
$(\rev(V_2),\rev(U_2))$.  With the same random $\Delta$ in both directions,
\[
 \mathbb E_\Delta\left[
 \LCS(\Trim_\Delta(U_1U_2),\Trim_{-\Delta}(V_1V_2))+\Delta
 \right]
 \ge2B+0.08\eps B.
\]
\end{lemma}

\begin{proof}
The forward matching uses the first $S$ ones of $U_1$ and a $T$-one interval
of $V_1$ beginning after $\Delta$ ones.  The reversed matching maps to a
$T$-one interval of $U_2$ ending $\Delta$ ones before its end and the last
$S$ ones of $V_2$.  The middle regions in the two concatenations each contain
exactly $2B-S-T-\Delta$ ones.  Match those ones and apply
\eqref{eq:BY-reward} twice.  Linearity of expectation is sufficient; no
independence is used.
\end{proof}

For fixed selected source flags, the bad-weight function of $\Delta$ is a
sum, over the $O(\log B)$ powers $p$, of cross-correlations between a sparse
weighted source array and the target bad-start indicator at scale $Q(p)$.
Thus all values $W(\Delta)$ are computed deterministically in
$B\polylog B$ time.

\section{Exact stitching by diagonal total variation}
\label{sec:stitching}

Fix a baseline symbol $b\in\{0,1\}$.  A local certificate $e$ consists of
ordered source and target physical intervals, an explicit common subsequence
of weight $w(e)$, and baseline counts $q_s(e),q_t(e)$ in those intervals.
Let $R_b^x(u)$ denote the number of copies of $b$ strictly before physical
position $u$ in $x$.  Define
\[
 z^-(e)=R_b^t(\ell_t(e))-R_b^s(\ell_s(e)),
 \qquad
 z^+(e)=z^-(e)+q_t(e)-q_s(e),
\]
and
\[
 r(e)=w(e)-\frac{q_s(e)+q_t(e)}2.
\]

\begin{lemma}[Exact chain identity]\label{lem:chain}
Let $M_s,M_t$ be the total baseline counts in $s,t$.  For any ordered chain
$e_1\prec\cdots\prec e_k$, concatenating the local certificates and matching
the baseline symbol greedily in every gap produces a common subsequence of
length at least
\begin{align}\label{eq:chain}
 \frac{M_s+M_t}{2}+\sum_{i=1}^kr(e_i)
 -\frac12\Bigl(&|z^-(e_1)|
 +\sum_{i=1}^{k-1}|z^-(e_{i+1})-z^+(e_i)|\\
 &+|(M_t-M_s)-z^+(e_k)|\Bigr).
\end{align}
\end{lemma}

\begin{proof}
If a source gap contains $x$ baseline symbols and its target gap contains
$y$, greedy matching contributes
$\min(x,y)=(x+y-|x-y|)/2$.  Summing the $x+y$ terms over all gaps gives
$M_s+M_t-\sum_i(q_s(e_i)+q_t(e_i))$.  The gap differences are exactly the
three kinds of diagonal differences in \eqref{eq:chain}.  Adding the local
weights proves the identity.
\end{proof}

Consequently the optimal chain over a finite certificate set is found by
\begin{equation}\label{eq:DP}
 F(e)=r(e)+\max\left\{-\frac{|z^-(e)|}{2},
 \max_{f\prec e}\left(F(f)-\frac{|z^-(e)-z^+(f)|}{2}\right)\right\}.
\end{equation}
After computing these values, the length certified by the best chain is
\begin{equation}\label{eq:DP-terminal}
 \frac{M_s+M_t}{2}+
 \max\left\{0,\max_e\left(
 F(e)-\frac{|(M_t-M_s)-z^+(e)|}{2}\right)\right\}.
\end{equation}
Predecessor pointers recover the chain.  The zero in
\eqref{eq:DP-terminal} represents the empty chain.
There are only $\polylog N$ certificates at a fixed scale and shift, so a
quadratic implementation is sufficient.

\begin{lemma}[Hidden offset variation]\label{lem:TV}
For a chain arising from the fixed common subsequence $q$, the extended total
variation of its unshifted base diagonals is
\[
 O(D+Kh)=O(D+\sigma N).
\]
This assertion applies in either of the following two situations:
\begin{enumerate}[label=(\roman*)]
\item the certificates use full outward-rounded physical witness spans and
      the baseline is either symbol;
\item the certificates lie in first-$B$ or last-$B$ canonical windows and
      the baseline is one.
\end{enumerate}
\end{lemma}

\begin{proof}
At a cut of the fixed matching, the difference between the target and source
$b$-ranks equals the difference between the numbers of unmatched $b$ symbols
seen so far.  Along increasing cuts, the total variation of this quantity is
at most the combined number of unmatched $b$ symbols, at most $2D$.

In case (i), moving a true physical endpoint outward to its grid point changes
either symbol-rank by less than $h$.  In case (ii), the one-rank of a first or
last canonical boundary differs from the corresponding core one-rank by at
most the number of additional ones in that occurrence.  On the disjoint
parity, the sum of these additional-one errors is at most
$2D+4Kh$ by \Cref{lem:defect,lem:grid}.  The triangle inequality proves the
claim, including the initial and terminal diagonals.  Notice that the
canonical assertion would be false for the zero baseline, which is why
case (ii) is explicitly restricted to the one baseline.
\end{proof}

For Green certificates, one common residual shift $\Delta$ adds the same
quantity to the start and end diagonal of every selected block.  It therefore
costs only $|\Delta|$ at the two global endpoints.  For Blue--Yellow, a
forward certificate has
\[
 (z^-,z^+)=(P_i^F+\Delta,\ P_i^F+\Delta+(T-S)),
\]
while a later reversed certificate has
\[
 (z^-,z^+)=(P_j^R+\Delta+(T-S),\ P_j^R+\Delta).
\]
Thus the excursion $T-S$ cancels exactly; the intervening penalty is only
$|P_j^R-P_i^F|/2$, already covered by \Cref{lem:TV}.

\begin{corollary}[Separated directional blocks]\label{cor:BY-separated}
Consider any ordered alternating chain of forward and later reverse
Blue--Yellow certificates, all using the same $\Delta\in[0,B/4]$.  Relative
to the all-one baseline, the chain's adjusted reward is the sum of its local
adjusted rewards, minus $O(D+Kh+B)$.  In particular the conclusion uses no
adjacency assumption on paired directional blocks.
\end{corollary}

\begin{proof}
For a forward block $i$ the two diagonals are
$(P_i^F+\Delta,P_i^F+\Delta+T-S)$; for a later reverse block $j$ they are
$(P_j^R+\Delta+T-S,P_j^R+\Delta)$.  Thus each $T-S$ excursion cancels at the
next forward--reverse interface.  All remaining internal penalties are the
variation of the unshifted $P$ values and are $O(D+Kh)$ by
\Cref{lem:TV}.  The unmatched excursion, if the chain has odd length, and
the two terminal shifts cost $O(B)$ in total.  Apply \Cref{lem:chain}.
\end{proof}

\section{The balanced algorithm and its running time}
\label{sec:balanced-algorithm}

Choose an absolute exponent
$A>\max\{A_0,200\gamma^{-3}\}+10$ and a sufficiently large absolute
constant $K_0$; $K_0\ge10^6\eps^{-6}$ is enough.  For all powers of two
\begin{equation}\label{eq:B-range}
 \frac{N}{(\log N)^A}\le B\le\frac{N}{K_0},
\end{equation}
the algorithm performs the following operations.  Reverse-oriented local
tests use $\rev(L_B(X))$; the full inputs are never passed to the modified
reversal.

\begin{enumerate}[label=\arabic*.]
\item Set $h=\lceil\sigma B\rceil$ and enumerate every physical interval
      with endpoints on the $h$-grid in each input.
\item For every physical interval pair, store both unary certificates:
      matching as many zeros as possible and matching as many ones as
      possible.
\item For every physical interval containing at least $B$ ones, form its
      first-$B$ and last-$B$ canonical variants.
\item For every pair of canonical variants, construct the weak Green flag
      arrays at every scale $r=g+1$, $g\in\mathcal G_B$, compute all shift
      correlations, and store the certified Green score functions.  Do this
      both left-to-right on first canonical windows and right-to-left on
      reversed last canonical windows.
\item For every canonical pair and every integer
      $m_0\in[1,\lfloor\log_2 B\rfloor]$, test the two observable hypotheses
      of \Cref{lem:BY-directional}.  Do this in both ordered directions:
      $(F_B(X_s),F_B(X_t))$ and
      $(\rev(L_B(X_t)),\rev(L_B(X_s)))$.  When the hypotheses hold, greedily
      select source flags and compute all directional Blue--Yellow score
      functions by grouped cross-correlations.
\item For every residual shift $\Delta$ in the appropriate range and for
      each baseline symbol, run \eqref{eq:DP} on the resulting ordered
      certificates, evaluate \eqref{eq:DP-terminal}, and reconstruct the best
      explicit chain.  Compare it with the two whole-input unary
      subsequences.
\end{enumerate}

For $N$ below the absolute threshold needed by the structural lemmas, use the
quadratic exact LCS algorithm.  This does not affect the asymptotic bound,
because the threshold is an absolute constant.

\begin{lemma}[Soundness and running time]\label{lem:runtime}
Every reported value is the length of an explicit common subsequence.  The
algorithm is deterministic and runs in $N\polylog N$ time on a word RAM with
$\Theta(\log N)$-bit words.
\end{lemma}

\begin{proof}
Soundness follows from the explicit local constructions and
\Cref{lem:chain}.  Prefix counts and select data structures give constant-time
symbol-count and one-rank interval queries.

At a scale in \eqref{eq:B-range}, the physical grid has
$G=O(N/B)=\polylog N$ points.  There are $O(G^2)$ intervals per input and
$O(G^4)$ interval pairs.  Green and Blue--Yellow processing costs
$B\polylog B$ per pair, so its total is
\[
 O(G^4B\polylog B)=N\polylog N.
\]
There are $O(\log\log N)$ candidate scales.  At a fixed $\Delta$, the
certificate set has size $\polylog N$; running its quadratic DP for every one
of $O(B)$ shifts costs $B\polylog N=N\polylog N$.  All correlation inputs
are nonnegative integer arrays of length $O(B)$ with entries bounded by a
fixed polynomial in $N$.  Choose a base $2^w$, $w=O(\log N)$, larger than
twice the largest possible convolution coefficient and pack the two
polynomials as integers (reversing one coefficient sequence for
cross-correlation).  Kronecker substitution then recovers every coefficient
from one exact product of two $O(B\log N)$-bit integers.  The deterministic
Sch\"onhage--Strassen multiplication algorithm~\cite{SchonhageStrassen}
therefore gives $B\polylog N$ bit operations, and hence at most that many
word operations.  No numerical approximation or randomized prime choice is
involved.  Reconstruction is linear in the total selected span.
\end{proof}

\section{Completeness under a near-perfect common subsequence}
\label{sec:balanced-completeness}

\begin{lemma}[The successful scale is tried]\label{lem:scale-tried}
Assume $D\le\eta N$.  For all sufficiently large $N$, the scale $B_0$ in
\Cref{prop:GHL}(a), and every scale $B$ furnished by
\Cref{prop:GHL}(b) after choosing
$R=\lceil\log_2K_0\rceil$, lie in the range \eqref{eq:B-range}.
\end{lemma}

\begin{proof}
In alternative (a), the floor in the definition of $B_0$ gives
\[
 \frac{L}{2(\log_2L)^{200\gamma^{-3}}}
 <B_0\le \frac{L}{(\log_2L)^{200\gamma^{-3}}}.
\]
In alternative (b), $K\le(\log L)^{A_0}$ gives
$B=L/K\ge L/(\log L)^{A_0}$, while the choice of $R$ gives
$K\ge K_0$.  Now use $L>N/4-D/2$, $L\le N/2$, and the choice of $A$.
The upper bound $B_0\le N/K_0$ also holds for all sufficiently large $N$.
\end{proof}

Assume $D\le\eta N$ and use the hidden $q'$ and $L$ from
\eqref{eq:L-linear}.  The algorithm does not know them; they are used only to
identify a successful candidate chain among the enumerated objects.
\Cref{lem:scale-tried} ensures that the required grid scale is enumerated.

First suppose Proposition~\ref{prop:GHL}(a) holds.  Round the extracted
zero-rich interval occurrences on the physical grid at scale $B_0$ and take
the heavier parity.  Their total one-rank size remains at least
$\gamma^2L/4$.  Indeed, order the extracted intervals.  Between intervals
$j$ and $j+2$ lies interval $j+1$, whose two physical occurrence spans each
have length at least $B_0$.  Since $2h<B_0$, the outward-rounded spans of
either fixed parity are disjoint in both inputs; weighted parity
pigeonholing preserves at least half the total one-rank mass.  If there are
$M$ selected intervals, then
$M\le L/B_0$ before parity selection.  Their total additional-symbol count
after rounding is $O(D+Mh)=O(D+\sigma L)$ by defect charging and
\eqref{eq:Kh}.  Summing \Cref{lem:imbalanced}(a) over the selected intervals
and using the global all-one baseline therefore gives
\begin{equation}\label{eq:huge-gain}
 \operatorname{ALG}(s,t)
 \ge\frac N2+\Omega(\eps\gamma^2L)-O(D+\sigma N)
 =\frac N2+\Omega(\eps^6L)-O(D+\sigma N).
\end{equation}

We may therefore assume Proposition~\ref{prop:GHL}(b).  Choose its scale with
$R$ large enough that $K=L/B\ge K_0$.  A block is \emph{defective} if
$e_i+4h>\kappa B$.  By \Cref{lem:defect}, the defective fraction is
\[
 O(D/(\kappa L)+\sigma/\kappa+1/(\kappa B))=o(1).
\]
All nondefective rounded occurrences satisfy the hypotheses of
\Cref{lem:canonical-stability}.

Assign priority types independently to every block and its reversal.  We
consider three exhaustive cases.

\paragraph{Case 1: many Imbalanced blocks.}
Suppose at least $K/10$ blocks are assigned Imbalanced in one orientation.
Discard defective blocks, choose a parity with at least half the survivors,
and pigeonhole the sign of imbalance.  At least $K/50$ blocks remain for
small enough $\eta,\sigma$.  Every witness has one-rank size
$r>\eps^5B$.  Use the all-one baseline for zero-rich witnesses and the
all-zero baseline for zero-poor witnesses; because the sign was
pigeonholed, one single baseline is used by the entire chain.  By
\Cref{lem:imbalanced}, every
selected block has adjusted reward at least $\eps^6B/3$.  The chain identity
and \Cref{lem:TV} yield
\begin{equation}\label{eq:I-gain}
 \operatorname{ALG}(s,t)
 \ge\frac N2+\frac{\eps^6}{150}L-O(D+\sigma N).
\end{equation}

\paragraph{Case 2: many Green blocks.}
Suppose Case 1 does not hold and at least $K/10$ blocks are assigned Green in
one orientation.  After defect removal and parity selection, at least $K/30$
remain.  Their Green scales may differ.  By
\Cref{lem:canonical-stability,lem:green}, each has expected adjusted reward
at least $0.05\eps^5B$ under the same uniform residual shift
$\Delta\in[-B,B]$.  Linearity of expectation gives one common $\Delta$ for
which the sum of local rewards has this lower bound.  The residual shift is
paid only once, so
\begin{equation}\label{eq:G-gain}
 \operatorname{ALG}(s,t)
 \ge\frac N2+\frac{\eps^5}{600}L-O(D+\sigma N+B).
\end{equation}

\paragraph{Case 3: Blue--Yellow in both orientations.}
If neither of the first two cases occurs in either orientation, more than
$4K/5$ blocks are assigned Blue--Yellow forward and more than $4K/5$ are
assigned Blue--Yellow after reversal.  Intersect these sets with the two
regular sets and the nondefective set, then choose a parity.  At least $K/4$
survivors remain for large enough $N$.  Order them and use them alternately
as forward and reverse blocks, producing at least $K/10$ directional pairs.

By \Cref{lem:canonical-stability,lem:BY-directional}, each direction has
expected adjusted reward at least $0.04\eps B$ for one common
$\Delta\in[0,B/4]$.  The forward--reverse excursions cancel, and
\Cref{cor:BY-separated} controls the separated blocks.  Hence some common
$\Delta$ satisfies
\begin{equation}\label{eq:BY-gain}
 \operatorname{ALG}(s,t)
 \ge\frac N2+0.008\eps L-O(D+\sigma N+B).
\end{equation}

When a dense Imbalanced set occurs in the reverse orientation, map each
one-rank witness interval back to its original occurrence; symbol counts and
order are preserved.  When a dense Green set occurs in the reverse
orientation, use the right-to-left arrays of $\rev(L_B(X))$ and map the
resulting local matching back.  Thus neither argument requires a reversal of
an arbitrary full input.

Since $K\ge K_0$, the one-time $B$ loss in
\eqref{eq:G-gain} and \eqref{eq:BY-gain} is smaller than a fixed small fraction of the
corresponding gain.  Combining \eqref{eq:huge-gain}--\eqref{eq:BY-gain},
there is an absolute $a_*>0$ such that
\begin{equation}\label{eq:master}
 \operatorname{ALG}(s,t)
 \ge\frac N2+a_*L-C(D+\sigma N)
\end{equation}
for an absolute $C$.  By \eqref{eq:L-linear},
\[
 a_*L\ge a_*\left(N/4-D/2\right).
\]
Choose $\sigma$ and then $\eta$ sufficiently small relative to $a_*/C$.
For $D\le\eta N$, \eqref{eq:master} is at least
$(1/2+\delta)N$ for some absolute $\delta>0$.  This proves
\Cref{thm:balanced}.

\section{Reductions to arbitrary inputs}
\label{sec:reductions}

Let $\eta_0,\delta_0>0$ be constants for which
\Cref{thm:balanced} holds.  We may assume $\delta_0\le1/2$.  Fix
\begin{equation}\label{eq:near-balance-r}
 r=\min\left\{\frac1{40},\frac{\eta_0}{20},
                    \frac{\delta_0}{11}\right\}.
\end{equation}
For arbitrary binary strings $x,y$, write
\begin{equation}\label{eq:best-match}
 \BestMatch(x,y)=
 \max_{a\in\{0,1\}}\min\{\#_a(x),\#_a(y)\}.
\end{equation}
This is the length of the longer of the two unary common
subsequences and is computable, with its matching, in linear time.

\subsection{Nearly balanced equal-length strings}

Call a length-$m$ string $r$-\emph{nearly balanced} if each of its
symbol counts lies in
$[(1/2-r)m,(1/2+r)m]$.

\begin{lemma}[Trimming lemma]\label{lem:nearly-balanced}
There is a deterministic $m\polylog(m+2)$-time algorithm which, on
two $r$-nearly-balanced strings $x,y\in\{0,1\}^m$, outputs an explicit
common subsequence of length at least
\[
 \left(\frac12+r\right)\LCS(x,y).
\]
\end{lemma}

\begin{proof}
If $m=0$, return the empty subsequence.  Henceforth $m>0$.
Put
\[
 k=\min\{\#_0(x),\#_1(x),\#_0(y),\#_1(y)\},
 \qquad M=2k.
\]
Retain exactly $k$ occurrences of each symbol in each string, and let
$x',y'$ be the resulting strings.  They are exactly balanced and have
common length $M$.  Near balance gives
\begin{equation}\label{eq:M-lower}
 M\ge(1-2r)m,
 \qquad m-M\le2rm.
\end{equation}

Let $L=\LCS(x,y)$ and $L'=\LCS(x',y')$.  Fix an optimal matching for
$x,y$.  Deleting $m-M$ positions from each input destroys at most the
sum of those two numbers of matched pairs.  Thus
\begin{equation}\label{eq:trim-loss}
 L'\ge L-2(m-M).
\end{equation}
The factor two is important: the deletions made in the two inputs need
not destroy the same matched pairs.

The algorithm returns the better of $\BestMatch(x,y)$ and the output
of the balanced promise algorithm on $x',y'$.  Since every symbol
count in both inputs is at least $k$,
$\BestMatch(x,y)\ge k=M/2$.  If this candidate is shorter than
$(1/2+r)L$, then
\begin{equation}\label{eq:L-near-M}
 L>\frac{M}{1+2r}.
\end{equation}
Combining \eqref{eq:trim-loss}, \eqref{eq:L-near-M}, and
\eqref{eq:M-lower},
\begin{align}
 M-L'
 &\le M-L+2(m-M)\nonumber\\
 &<\frac{2r}{1+2r}M+2(m-M)
 \le6rm.\label{eq:trim-defect}
\end{align}
Moreover,
\[
 \frac{6rm}{M}\le\frac{6r}{1-2r}<\eta_0
\]
by \eqref{eq:near-balance-r}.  Hence the promise in
\Cref{thm:balanced} holds for $x',y'$, and that algorithm returns at
least $(1/2+\delta_0)M$ symbols.  Finally,
\begin{align*}
 \left(\frac12+\delta_0\right)M
 &\ge\left(\frac12+\delta_0\right)(1-2r)m\\
 &\ge\left(\frac12+r\right)m
 \ge\left(\frac12+r\right)L.
\end{align*}
Indeed, the difference between the second line's two coefficients is
$\delta_0-2r(1+\delta_0)>0$ by
$r\le\delta_0/11$ and $\delta_0\le1/2$.

One pass through each input performs the trimming while retaining a
map to the original positions.  The balanced matching therefore lifts
directly to $x,y$.  This pass and reconstruction are linear, and the
only nontrivial call takes $M\polylog(M+2)$ time.
\end{proof}

\subsection{Arbitrary equal-length strings}

We next modify the equal-length reduction of Rubinstein and
Song~\cite{RubinsteinSong}.  Their proof normalizes all lengths and
symbol counts by the common input length.  We retain that convention
in this subsection.  Thus $\#_a(u)$ below means the number of
$a$'s in $u$ divided by the original length $n$, and $\LCS(u,v)$ is
normalized by $n$ as well.

For ordered strings $U,V,W$, let
\[
 \Greedy(U,V;W)
 =\max_{W=W_1W_2}
   \bigl(\BestMatch(U,W_1)+\BestMatch(V,W_2)\bigr),
\]
where the maximum is over contiguous prefix--suffix decompositions.
A single sweep computes an optimizing split and the associated
matching.  The notation $\Match(U,V,a)$ denotes the unary matching of
symbol $a$ between $U$ and $V$.

\begin{theorem}[Equal-length reduction]\label{thm:equal-length}
There is an absolute constant $\delta_{\rm eq}>0$ and a deterministic
$n\polylog(n+2)$-time algorithm which, for arbitrary
$A,B\in\{0,1\}^n$, outputs an explicit common subsequence of length at
least
\[
 \left(\frac12+\delta_{\rm eq}\right)\LCS(A,B).
\]
\end{theorem}

\begin{proof}
If $n$ is below a fixed absolute threshold, compute an exact LCS by dynamic
programming and return it.  Thus $n>0$, so the normalized quantities used in
this subsection are well defined.
Set
\begin{equation}\label{eq:RS-parameters}
 b=\frac r{100},\qquad \tau=\frac b2,
 \qquad \delta_{\rm eq}=\frac b{10}.
\end{equation}
Run $\BestMatch(A,B)$ at the outset.  By interchanging the inputs and
complementing both strings if necessary, assume that
\[
 \alpha=\min\{\#_0(A),\#_1(A),\#_0(B),\#_1(B)\}
 =\#_1(A).
\]
These symmetries preserve both LCS and explicit matchings.  If
$\alpha=0$, $A$ has only one symbol and $\BestMatch$ is optimal, so
assume $\alpha>0$.

We first dispose of the cross-unbalanced case.  If
\begin{equation}\label{eq:cross-unbalanced}
 |\#_1(A)-\#_0(B)|>\tau\alpha,
\end{equation}
then minimality of $\alpha$ implies
$\#_0(B)>(1+\tau)\alpha$.  Also
$\BestMatch(A,B)=\#_0(B)$ and
\[
 \LCS(A,B)\le\alpha+\#_0(B).
\]
Consequently
\begin{equation}\label{eq:cross-ratio}
 \frac{\BestMatch(A,B)}{\LCS(A,B)}
 \ge\frac{1+\tau}{2+\tau}
 =\frac12+\frac{\tau}{4+2\tau}
 \ge\frac12+\delta_{\rm eq}.
\end{equation}
Thus the cross-count argument contributes the gain
$\tau/(4+2\tau)$ used below.

We may now assume
\begin{equation}\label{eq:cross-close}
 \alpha\le\#_0(B)\le(1+\tau)\alpha.
\end{equation}
If $\alpha>1/2-10b\alpha$, then every symbol count of both full
inputs is within
$(10b+\tau)\alpha<r$ of $1/2$.  The pair is therefore $r$-nearly
balanced, and \Cref{lem:nearly-balanced} applies.  Henceforth assume
\begin{equation}\label{eq:alpha-away-half}
 \alpha\le\frac12-10\beta,
 \qquad \beta=b\alpha.
\end{equation}

Partition both inputs into three contiguous pieces,
\[
 A=L_AM_AR_A,\qquad B=L_BM_BR_B,
\]
with all four end pieces of normalized length $\alpha$.  Put
\[
 a_L=\#_0(L_A),\quad a_R=\#_0(R_A),\qquad
 b_L=\#_1(L_B),\quad b_R=\#_1(R_B).
\]
The following six cases are exhaustive:
\begin{enumerate}[label=\textup{(\arabic*)},leftmargin=2.4em]
\item $b_R,a_R\le\alpha/2+2\beta$;
\item $b_L,a_L\le\alpha/2+2\beta$;
\item $b_L,b_R\le\alpha/2+\beta$ and
      $a_L,a_R>\alpha/2+2\beta$;
\item $b_L,b_R>\alpha/2+2\beta$ and
      $a_L,a_R\le\alpha/2+\beta$;
\item $b_R,a_L>\alpha/2+\beta$;
\item $b_L,a_R>\alpha/2+\beta$.
\end{enumerate}
For completeness, suppose neither (5) nor (6) holds.  If (1) fails
because $b_R>\alpha/2+2\beta$, then $a_L\le\alpha/2+\beta$;
failure of (2) forces $b_L>\alpha/2+2\beta$, which in turn forces
$a_R\le\alpha/2+\beta$, giving (4).  If instead (1) fails because
$a_R>\alpha/2+2\beta$, the symmetric argument gives (3).
Thus failure of (1), (2), (5), and (6) implies (3) or (4).

All six conditions depend only on observed symbol counts.  The
algorithm computes every candidate prescribed by every condition that
holds, including the reversed copies below, and returns the longest
candidate together with the initial $\BestMatch(A,B)$ output.

We verify the certificate in every case.  In Case (1), first suppose
\begin{equation}\label{eq:case-1a}
 a_R,b_R\in[\alpha/2-4\beta,\alpha/2+4\beta].
\end{equation}
The pair $(R_A,R_B)$ is $4b$-nearly balanced relative to its common
length $\alpha n$, and hence is within the range of
\Cref{lem:nearly-balanced} because $4b<r$.

The algorithm computes all three candidates
\begin{align}
 G_{AB}&=\Greedy(L_AM_A,R_A;B),\nonumber\\
 G_{BA}&=\Greedy(L_BM_B,R_B;A),\label{eq:three-candidates}\\
 P&=\BestMatch(L_AM_A,L_BM_B)
       +\operatorname{NB}(R_A,R_B),\nonumber
\end{align}
where $\operatorname{NB}$ is the algorithm from
\Cref{lem:nearly-balanced}.  All three candidates are needed because
the proof may branch on a partition induced by an unknown optimal
matching, while the algorithm has access only to the explicitly
computed greedy splits.

Fix an optimal matching.  If it has no match from $R_A$ into
$L_BM_B$, it induces a decomposition $B=\widehat L_B\widehat R_B$
for which
\[
 \LCS(A,B)
 =\LCS(L_AM_A,\widehat L_B)
  +\LCS(R_A,\widehat R_B).
\]
Otherwise, noncrossing implies that it has no match from $L_AM_A$
into $R_B$, and the transposed decomposition used by $G_{BA}$ applies.
Consider the applicable direction.  Let $X,Y$ be the symbol-count
upper bounds for the left and right LCS terms, and let $Z$ be the
larger of the two unary contributions to $X$.  Then
\begin{equation}\label{eq:XYZ}
 \LCS(A,B)\le X+Y,\qquad Z\ge X/2,\qquad
 G_{AB}\text{ or }G_{BA}\ge Z+Y/2.
\end{equation}
Condition \eqref{eq:case-1a} and \eqref{eq:cross-close} give
$X-Z\le\alpha/2+4\beta+\tau\alpha$.  Indeed, in the $G_{AB}$
direction the smaller unary contribution to $X$ is at most
$\#_1(L_AM_A)=a_R\le\alpha/2+4\beta$.  In the $G_{BA}$ direction it
is at most
\[
 \#_0(L_BM_B)=\#_0(B)-\alpha+b_R
 \le\alpha/2+4\beta+\tau\alpha.
\]
Hence, if
$Z>\alpha/2+10\beta$, the applicable greedy candidate is at least
\[
 \frac12\LCS(A,B)+2\beta.
\]

If $Z\le\alpha/2+10\beta$, then
\begin{equation}\label{eq:case-1a-upper}
 \LCS(A,B)\le\alpha+20\beta+\LCS(R_A,R_B).
\end{equation}
Furthermore
\[
 \BestMatch(L_AM_A,L_BM_B)\ge\alpha/2-4\beta
\]
and
$\LCS(R_A,R_B)\ge\alpha/2-4\beta$.  Thus the candidate $P$ exceeds
half the right side of \eqref{eq:case-1a-upper} by at least
\begin{align*}
 r\LCS(R_A,R_B)-14\beta
 &\ge\bigl(r(1/2-4b)-14b\bigr)\alpha
 \ge30\beta.
\end{align*}
Here we used $b=r/100$ and $r\le1/40$.

It remains in Case (1) to treat the complement of
\eqref{eq:case-1a}.  Then either
\[
 b_R<\alpha/2-4\beta,\quad a_R\le\alpha/2+2\beta,
\]
or the same statement holds with $a_R,b_R$ interchanged.  In either
alternative,
\begin{equation}\label{eq:right-count-sum}
 a_R+b_R\le\alpha-2\beta.
\end{equation}
Splitting an optimal noncrossing matching at the right blocks, the two
possible matching directions give the bounds
\[
 a_R+\#_0(B)+b_R
 \le2\alpha-2\beta+\tau\alpha
\]
and
\[
 (\#_0(B)-\alpha+b_R)+\alpha+a_R
 \le2\alpha-2\beta+\tau\alpha.
\]
The latter charges zeros in the left $B$-piece, all ones of $A$, and
zeros in the right $A$-piece; these three classes cover both matching
subproblems.
The two displayed expressions are equal to
$\#_0(B)+a_R+b_R$.  Thus \eqref{eq:right-count-sum} and
\eqref{eq:cross-close} prove, in both alternatives,
\begin{equation}\label{eq:case-1bc-upper}
 \LCS(A,B)\le2\alpha-2\beta+\tau\alpha.
\end{equation}
Since
$\BestMatch(A,B)\ge\alpha$, that upper bound gives an additive advantage over
half the LCS of at least
\begin{equation}\label{eq:case-1bc-gain}
 \beta-\frac{\tau\alpha}{2}=\frac34\beta.
\end{equation}
Case (2) is obtained by reversing both strings; the algorithm runs the
corresponding candidates.

In Case (3), concatenate
\begin{equation}\label{eq:case-3-certificate}
 \Match(L_A,L_B,0),\quad \Match(M_A,M_B,1),\quad
 \Match(R_A,R_B,0).
\end{equation}
The two outer terms each have length at least
$\alpha/2-\beta$.  The middle of $A$ has more than $4\beta$ ones.
The middle of $B$ has at least
\[
 1-(1+\tau)\alpha-(\alpha+2\beta)>4\beta
\]
ones by \eqref{eq:alpha-away-half}.  Thus
\eqref{eq:case-3-certificate} has length at least
$\alpha+2\beta$.  In Case (4), concatenate the one matchings on the
two end pairs and the zero matching on the middle pair.  Each outer
term has length at least $\alpha/2-\beta$.  The middle of $B$ has
\[
 \#_0(M_B)=\#_0(B)-(\alpha-b_L)-(\alpha-b_R)>4\beta,
\]
where we used $\#_0(B)\ge\alpha$ and the case assumptions.  The
middle of $A$ has
\[
 \#_0(M_A)=1-\alpha-a_L-a_R
 \ge1-2\alpha-2\beta>4\beta
\]
by \eqref{eq:alpha-away-half}.  This gives the same total lower bound
$\alpha+2\beta$ without invoking an exact symmetry.

In Case (5), concatenate
\begin{equation}\label{eq:case-5-certificate}
 \Match(L_A,L_BM_B,0),\qquad
 \Match(M_AR_A,R_B,1).
\end{equation}
Each term has length greater than $\alpha/2+\beta$: for instance,
\[
 \#_0(L_BM_B)=\#_0(B)-(\alpha-b_R)
 \ge b_R>\alpha/2+\beta,
\]
and the other three required counts follow directly from the case
assumptions.  Hence \eqref{eq:case-5-certificate} again has length
greater than $\alpha+2\beta$.  Case (6) follows by reversal.

Finally, throughout the current branch,
\begin{equation}\label{eq:L-alpha-upper}
 \LCS(A,B)\le
 \min\{\#_0(A),\#_0(B)\}+\min\{\#_1(A),\#_1(B)\}
 \le(2+\tau)\alpha.
\end{equation}
The smallest additive gain established above is $3\beta/4$.
Together with \eqref{eq:RS-parameters} and
\eqref{eq:L-alpha-upper}, this is more than
$\delta_{\rm eq}\LCS(A,B)$.  The cross-unbalanced and nearly balanced
branches already have the same guarantee.  This completes the case
analysis.

Every certificate used above is unary or is returned explicitly by
\Cref{lem:nearly-balanced}; concatenated certificates lie in ordered,
disjoint subproblems.  There are only constantly many oracle calls,
and all counts, greedy splits, symmetries, and reconstruction maps take
linear time.  The total running time is therefore
$n\polylog(n+2)$.
\end{proof}

Only the following uses of $\operatorname{ApproxED}$ from
Rubinstein--Song~\cite{RubinsteinSong} enter the preceding proof: its
approximately balanced Lemma 3.2, the Case 1(a) call on $(R_A,R_B)$,
and the reversed copy of that case.
All other cases are precisely the unary and greedy constructions
verified above.  Thus no edit-distance guarantee remains hidden in
\Cref{thm:equal-length}.

\subsection{Unequal lengths and the final factor}

We use the following quantitative form of the He--Li reduction.  It is
stated in the discussion immediately following their main theorem and
is implemented by their construction~\cite{HeLi}.

\begin{proposition}[He--Li equal-to-unequal reduction]
\label{prop:HeLi}
There is an absolute constant $A\ge1$ with the following property.
Suppose that, for every $k$, a deterministic algorithm on
equal-length binary strings of length at most $k$ runs in
$k\polylog(k+2)$ time and outputs an explicit common subsequence of
length at least $(1/2+\zeta)$ times optimal.  Then arbitrary binary
strings $x,y$ of maximum length $N$ admit a deterministic explicit
$(1/2+\zeta^A)$-approximation in $N\polylog(N+2)$ time.
\end{proposition}

We record the bookkeeping bounds needed to apply the He--Li
construction in this form.  They also determine how its hierarchy of
constants is chosen.

First consider a near-equal oracle call on lengths $p\le q$, where
$q-p\le\lambda p$ and completeness supplies an optimum
$L\ge(1-\chi)p$.  Delete $q-p$ symbols from the longer local string.
The equal-length instance has optimum $L'$ satisfying
\begin{equation}\label{eq:near-equal-loss}
 L'\ge L-(q-p)
 \ge\left(1-\frac{\lambda}{1-\chi}\right)L.
\end{equation}
Choose the scale parameters so that
\begin{equation}\label{eq:near-equal-hierarchy}
 \frac{\lambda}{1-\chi}\le\frac\zeta2.
\end{equation}
Then an equal-length $(1/2+\zeta)$-approximation on the trimmed pair
returns at least $(1/2+\zeta/2)L$ symbols, because
$\zeta/2\le\zeta/(1+2\zeta)$.  He and Li choose $\lambda$ and $\chi$
as fixed powers of $\zeta$, so this requirement only changes absolute
exponents in their hierarchy.

Second, write
\[
 u=\min\{\#_0(x),\#_0(y)\},\qquad
 v=\min\{\#_1(x),\#_1(y)\},
\]
and suppose $u\ge v$.  If $u\ge(1+\theta)v$, the unary candidate has
ratio at least
\[
 \frac{1+\theta}{2+\theta}>\frac12+\frac\theta5.
\]
Otherwise $d=u-v<\theta v\le\theta L$, where
$L=\LCS(x,y)$, because the $v$ copies of the minority symbol form a
common subsequence.  A balancing step may delete $d$ symbols from
each input, so the safe bound is
\begin{equation}\label{eq:HeLi-two-sided}
 L'\ge L-2d\ge(1-2\theta)L.
\end{equation}
If the internal algorithm has gain $\Delta$, then for
$\theta\le\Delta/10$ its output on the trimmed instance is at least
\[
 \left(\frac12+\Delta\right)(1-2\theta)L
 \ge\left(\frac12+\frac{7\Delta}{10}\right)L,
\]
using $\Delta\le1/2$.

Finally, an earlier one-input balancing step assumes
$\#_1(x)\ge(1/2-\rho)|x|$.  The number deleted there is at most
$2\rho|x|$, rather than $\rho|x|$, while
$L\ge(1/2-\rho)|x|$.  Thus its surviving optimum obeys
\begin{equation}\label{eq:HeLi-one-sided}
 L'\ge
 \left(1-\frac{4\rho}{1-2\rho}\right)L.
\end{equation}
Taking $\rho\le\Delta/20$ in \eqref{eq:HeLi-one-sided} preserves at
least half of the internal gain.  The hierarchy in the He--Li
construction takes all three tolerances to be fixed powers of
$\zeta$; after the displayed losses, increasing the absolute exponent
$A$ preserves the form of \Cref{prop:HeLi}.

One oracle call is implicit in the imbalanced-input reduction used by
He and Li: the balanced end-block subcase of their cited linear-time
lemma invokes an equal-length approximate-edit-distance routine.  In
the present instantiation, replace that invocation by the assumed
equal-length LCS algorithm and retain its explicit certificate.  At a
block length $m$, the construction makes
$O((N/m)\polylog(N+2))$ calls on strings of length
$O(m/\log(m+2))$.  Since every such call costs its input length times
a polylogarithmic factor, the sum over all calls and all logarithmically
many scales is $N\polylog(N+2)$.

For explicit output, store the certificate attached to every accepted
rectangle and the backpointer of every dynamic-programming state.
Unary and structural certificates are embedded greedily, and deletion
position maps lift all oracle certificates to the original strings.
Fix lexicographically first choices in ties, so determinism is
preserved.  Inputs below the absolute block-scale threshold, including
empty strings, are handled by exact dynamic programming.  This
completes the verification of \Cref{prop:HeLi} for the stated
near-linear oracle.

\begin{proof}[Proof of \Cref{thm:main}]
If $|s|+|t|$ is below the absolute threshold in \Cref{prop:HeLi}, compute an
exact LCS and return it.  Henceforth $N=\max\{|s|,|t|\}>0$.
Apply \Cref{prop:HeLi} to the algorithm of
\Cref{thm:equal-length}, with
\[
 \zeta=\delta_{\rm eq},\qquad
 \delta_{\rm all}=\delta_{\rm eq}^{A}.
\]
Its running time is $T(N)=N\polylog(N+2)$, so the resulting running
time is still $N\polylog(N+2)$.  Since
\[
 N=\max\{|s|,|t|\}\le |s|+|t|\le2N,
\]
this is also $n\polylog(n+2)$ for $n=|s|+|t|$.

The returned subsequence has length at least
$(1/2+\delta_{\rm all})\LCS(s,t)$.  Its conventional approximation
factor is therefore at most
\begin{align*}
 \frac1{1/2+\delta_{\rm all}}
 &=\frac2{1+2\delta_{\rm all}}\\
 &=2-\frac{4\delta_{\rm all}}{1+2\delta_{\rm all}}.
\end{align*}
Taking
\[
 c=\frac{4\delta_{\rm eq}^{A}}
          {1+2\delta_{\rm eq}^{A}}>0
\]
proves the theorem.
\end{proof}

{\small
\bibliographystyle{alpha}
\bibliography{references}
}

\end{document}
